\documentclass[10pt,a4paper]{amsart}
\usepackage[margin=19mm]{geometry}
\usepackage{amsmath,amssymb,mathtools}
\usepackage[scr=rsfs,cal=euler]{mathalfa}
\usepackage{xfrac}
\usepackage[unicode,hidelinks,bookmarks=false]{hyperref}
\newcommand{\ie}{i.e.\ }
\newcommand{\cf}{cf.\ }
\newcommand{\resp}{respectively\ }
\newcommand{\Subsection}{\subsection}
\providecommand{\nobreakdash}{\nobreak\hbox{-}}
\setlength{\emergencystretch}{2em}
\multlinegap0pt
\newtheorem*{degreeprinciple}{Equivariant degree principle}
\newtheorem{thm}{Theorem}[section]
\newtheorem{cor}[thm]{Corollary}
\newtheorem{lemma}[thm]{Lemma}
\newcommand{\R}{\mathbb R}
\newcommand{\F}{\mathbb F}
\newcommand{\Z}{\mathbb Z}
\begin{document}
\title[Borsuk-Ulam type theorem for the orthogonal group and orthog]{Borsuk-Ulam type theorem for the orthogonal group and orthogonal four-partitions}
\author{Oleg R. Musin}
\date{}
\maketitle
\noindent\emph{Source and licence.} Borsuk-Ulam type theorem for the orthogonal group and orthogonal four-partitions. Version arXiv:2609.09688v3 (CC BY 4.0). This material is distributed under the Creative Commons Attribution 4.0 International licence, http://creativecommons.org/licenses/by/4.0/, as recorded in the arXiv OAI licence metadata for this exact version and shown on the abstract page https://arxiv.org/abs/2609.09688v3. Full rights notice: licenses/arxiv-2609.09688-CC-BY-4.0.txt.\par\smallskip
\noindent\textit{Editorial scope.} Counted unit: the original \texttt{thm} environment \texttt{main} at source lines 75--77 together with its complete proof at lines 140--216; the delivered document keeps source lines 44--217 so that every referenced label and every citation resolves inside it. Native editable source is the paper's own concatenated TeX; the AMS wrapper is editorial formatting only. No publisher class, upstream script or unrelated artwork is redistributed.
\section{Introduction}

The \emph{hyperoctahedral group} $B_k = A_k \rtimes S_k$ is the
semidirect product of the sign-change group
$A_k = (\mathbb{Z}/2)^k = \langle\lambda_1,\ldots,\lambda_k\rangle$
and the symmetric group $S_k$, where $S_k$ acts by
$\sigma(\lambda_i)=\lambda_{\sigma(i)}$.
The generator $\lambda_\ell \in A_k$ acts on $\mathbb{R}^k$ by
negating the $\ell$-th coordinate.

For each pair $\{i,j\}\subset[k]$ with $i<j$, let
$\mathbb{R}_{\lambda_i+\lambda_j}$ be the one-dimensional real
$A_k$-representation on which $\lambda_\ell$ acts by $-1$ if
$\ell\in\{i,j\}$ and by $+1$ otherwise.  Set
\[
 W_2^+ = \bigoplus_{1\le i<j\le k}\mathbb{R}_{\lambda_i+\lambda_j},
\]
extended to a $B_k$-representation by
$\sigma\,e_{\{i,j\}}=e_{\{\sigma(i),\sigma(j)\}}$.
One checks that $(W_2^+)^{B_k}=\{0\}$ and
$\dim W_2^+=\binom{k}{2}$.  Indeed, if a vector is fixed by
$\lambda_i$, every coefficient indexed by a pair containing $i$ must
equal its own negative; varying $i$ forces all coefficients to vanish.

\medskip

The Stiefel manifold $V_{k,k}=O(k)$ carries a free $B_k$-action:
$\lambda_\ell$ negates the $\ell$-th frame vector and $S_k$
permutes them.  Since $\dim W_2^+=\binom{k}{2}=\dim O(k)$,
the following zero theorem has the right dimensional setup.


\begin{thm}\label{main}
Every continuous $B_k$-equivariant map $f:O(k)\to W_2^+$ has a zero.
\end{thm}


\medskip

The classical ham-sandwich theorem~\cite{ST} guarantees that any
$d$ measures in $\mathbb{R}^d$ can be simultaneously bisected by a
single hyperplane.  For $d=2$, a single measure can already be
divided into four equal parts by two orthogonal lines~\cite{Shashkin}.

For general $d$, Makeev~\cite{Mak07} stated the following result.

\begin{thm}\label{mak}
Let $d\ge2$ and let $\mu$ be a finite Borel measure in $\mathbb{R}^d$
for which every hyperplane has measure zero.  Then there exist $d$
mutually orthogonal hyperplanes such that every pair of them divides
$\mu$ into four equal parts.
\end{thm}

Makeev outlines a dimension-count and parity argument, but leaves
without proof the key steps: uniqueness for the model example,
transversality, and the limiting argument.
Partial results were later obtained in~\cite{BK12,Mejia25}.
The present note gives a direct proof that supplies all three steps.

We deduce Theorem~\ref{mak} from Theorem~\ref{main} as follows.
For a measure $\mu$ with smooth positive density, each unit vector
$u$ determines a unique bisecting hyperplane, and the imbalance
of $\mu$ across each pair of hyperplanes defines a continuous
$B_d$-equivariant map $F:O(d)\to W_2^+$.
A zero of $F$ gives the desired orthogonal configuration;
the general case follows by Gaussian approximation.

The result is also a special case of the general orthogonal mass
partition theorem of~\cite{Mus26}, whose proof uses characteristic
classes on Stiefel manifolds.  The proof here instead uses a single
explicit model map and mod-$2$ degree.

\section{Proofs}

\subsection{Proof of Theorem~\ref{main}}\label{sec:proof-main}

We use the following standard mod-$2$ form of equivariant degree.

\begin{lemma}[Equivariant parity principle]\label{lem:degree}
Let a finite group $G$ act freely and smoothly on a closed
$n$-manifold $M$, and let $V$ be an $n$-dimensional real
$G$-representation.  Suppose that a smooth $G$-equivariant map
$h:M\to V$ has precisely one $G$-orbit of zeros and that every zero
is nondegenerate.  Then every continuous $G$-equivariant map
$f:M\to V$ has a zero.
\end{lemma}

\begin{proof}
Equivariant maps $M\to V$ are sections of the rank-$n$ bundle
$E=M\times_GV\to M/G$.  The section induced by $h$ has one transverse
zero, so
$\langle w_n(E),[M/G]\rangle=1\in\mathbb F_2$.
Thus $w_n(E)\ne0$, whereas a nowhere-zero section would force
$w_n(E)=0$.  The conclusion follows, with continuous sections handled
by approximation.  This is also the mod-$2$ equivariant degree
principle; see~\cite[Theorem~1]{Mus12}.
\end{proof}


\begin{proof}[Proof of Theorem~\ref{main}]
By Lemma~\ref{lem:degree}, it is enough to construct a smooth model map
whose zero set is a single free $B_k$-orbit and whose zeros are
nondegenerate.  We verify the three required facts:
\textit{Equivariance}, \textit{Zero set}, and \textit{Nondegeneracy}.

Fix pairwise distinct numbers $\alpha_1<\cdots<\alpha_k$ and, in the
standard orthonormal basis $e_1,\ldots,e_k$ of $\mathbb{R}^k$, take the
diagonal self-adjoint operator
\[
 A=\operatorname{diag}(\alpha_1,\ldots,\alpha_k),
 \qquad Ae_i=\alpha_i e_i.
\]
Define
\[
 h:O(k)\longrightarrow W_2^+,
 \qquad
 h(u_1,\ldots,u_k)
 =\sum_{1\le i<j\le k}
 \langle Au_i,u_j\rangle\, e_{\{i,j\}}.
\]

\textit{Equivariance.}
Negating $u_\ell$ negates $\langle Au_i,u_j\rangle$ exactly when
$\ell\in\{i,j\}$, matching the sign action on $e_{\{i,j\}}$.
Permuting the frame vectors permutes the index pairs.
Hence $h$ is $B_k$-equivariant.

\textit{Zero set.}
Write $U$ for the orthogonal matrix with columns $u_1,\ldots,u_k$.
The matrix of $A$ in this basis is $U^TAU$, and its $(i,j)$-entry is
\[
 (U^TAU)_{ij}=\langle Au_j,u_i\rangle.
\]
Since $A$ is self-adjoint, this matrix is symmetric and
$\langle Au_j,u_i\rangle=\langle Au_i,u_j\rangle$.  Consequently,
$h(U)=0$ exactly when every off-diagonal entry of $U^TAU$ vanishes,
that is, exactly when $U^TAU$ is diagonal.  In that case each column
$u_i$ is an eigenvector of $A$.  The spectrum is simple, so every
eigenspace is one-dimensional; hence
\[
 u_i=\varepsilon_i e_{\sigma(i)}
\]
for some signs $\varepsilon_i\in\{\pm1\}$ and a permutation
$\sigma\in S_k$.  Conversely, every signed permutation of the standard
basis is clearly a zero.  Therefore
$Z_h=B_k\cdot(e_1,\ldots,e_k)$ is a single free $B_k$-orbit.

\textit{Nondegeneracy.}
Identify tangent vectors to $O(k)$ at the eigenframe
$U_0=(e_1,\ldots,e_k)$ with skew-symmetric matrices $X=(x_{ij})$.
This follows by differentiating $U(t)^TU(t)=I$.  Along
$U(t)=U_0\exp(tX)$, the $i$-th column satisfies
\[
 u_i'(0)=\sum_p x_{pi}e_p.
\]
For $i<j$, the product rule gives
\begin{align*}
 (Dh)_{U_0}(X)_{\{i,j\}}
 &=\langle Au_i'(0),e_j\rangle
   +\langle Ae_i,u_j'(0)\rangle\\
 &=\alpha_jx_{ji}+\alpha_ix_{ij}
  =(\alpha_i-\alpha_j)x_{ij},
\end{align*}
because $x_{ji}=-x_{ij}$.  The variables $x_{ij}$, $i<j$, are
coordinates on $T_{U_0}O(k)$, and the same pairs index a basis of
$W_2^+$.  Thus $(Dh)_{U_0}$ is diagonal with nonzero diagonal entries
$\alpha_i-\alpha_j$.  With the coordinate pairs ordered in the same way,
\[
 \det (Dh)_{U_0}=\prod_{1\le i<j\le k}(\alpha_i-\alpha_j)\ne0.
\]
Hence $(Dh)_{U_0}$ is an isomorphism, so $U_0$ is an isolated zero and
is nondegenerate.  Equivariance gives the same conclusion at every
other zero.  Thus $Z_h$ is one free orbit of nondegenerate zeros, and
Lemma~\ref{lem:degree} shows that every continuous
$B_k$-equivariant map $O(k)\to W_2^+$ has a zero.
\end{proof}

\begin{thebibliography}{99}
\bibitem[BK12]{BK12} P.~Blagojevi\'{c} and R.~Karasev, Extensions of theorems of Rattray and Makeev, \textit{Topol.\ Methods Nonlinear Anal.} \textbf{40} (2012), no.\,1, 189--213.
\bibitem[Mak07]{Mak07} V.\,V.~Makeev, Equipartition of a continuous mass distribution, \textit{J.\ Math.\ Sci.} \textbf{140} (2007), no.\,4, 551--557.
\bibitem[Mejia25]{Mejia25} A.~Mej\'{\i}a, S.~Simon, and J.~Zhang, The generalized Makeev problem revisited, \textit{Beitr.\ Algebra Geom.} \textbf{66} (2025), 253--274.
\bibitem[Mus12]{Mus12} O.\,R.~Musin, Borsuk--Ulam type theorems for manifolds, \textit{Proc.\ Amer.\ Math.\ Soc.} \textbf{140} (2012), 2551--2560.
\bibitem[Mus26]{Mus26} O.\,R.~Musin, Borsuk--Ulam-type theorem for Stiefel manifolds and orthogonal mass partitions, preprint (2026), arXiv:2603.18550.
\bibitem[ST]{ST} A.\,H.~Stone and J.\,W.~Tukey, Generalized ``sandwich'' theorems, \textit{Duke Math.\ J.} \textbf{9} (1942), 356--359.
\bibitem[Shashkin]{Shashkin} Yu.\,A.~Shashkin, \textit{Fixed Points}, Amer.\ Math.\ Soc., Providence, RI, 1991.
\end{thebibliography}
\end{document}
